\section*{Hoofdstuk \ref{ch:b2:ism-star-formation} --- Het interstellaire medium en de vorming van sterren}

\begin{solution}{exo:b2:ism-star-formation:1}
$n = \num{2.7e-24}/\num{1.67e-24} = 1.6$ waterstofatomen per kubieke centimeter. In lucht is
$n = P/k_{\mathrm B}T = 101325/(\num{1.381e-23}\times288) = \qty{2.55e25}{m^{-3}}$,
dat is \num{2.5e19} moleculen per kubieke centimeter: ongeveer $1.6\times10^{19}$ keer
meer.
\end{solution}

\begin{solution}{exo:b2:ism-star-formation:2}
$M_J \propto n^{-1/2}$: $n$ delen door 100 vermenigvuldigt $M_J$ met 10, dus
$M_J = 54.6$ zonsmassa's bij $10^2$ deeltjes per kubieke centimeter. De wolk bevat
$\num{24000}/54.6 \approx 440$ \omterm{def:b2:ism-star-formation:jeans}{Jeans-massa}'s: ze kan als geheel niet stabiel zijn, en moet
fragmenteren.
\end{solution}

\begin{solution}{exo:b2:ism-star-formation:3}
$c_s = \sqrt{\num{1.381e-23}\times10/(2.33\times\num{1.661e-27})} =
\qty{189}{m.s^{-1}}$ bij \qty{10}{K}; $\sqrt{10}$ keer meer, \qty{597}{m.s^{-1}},
bij \qty{100}{K}. Doorkruisen van \qty{0.1}{pc}: $\num{3.086e15}/189 = \qty{1.63e13}{s}$,
dat is \qty{0.52}{Myr}.
\end{solution}

\begin{solution}{exo:b2:ism-star-formation:4}
$\rho = 2.33\times\num{1.661e-27}\times n$ met $n$ in \unit{m^{-3}}. Voor $10^5$
deeltjes per kubieke centimeter is $\rho = \qty{3.87e-16}{kg.m^{-3}}$ en
$t_{\mathrm{ff}} = \sqrt{3\pi/(32G\rho)} = \qty{3.38e12}{s} = \qty{0.107}{Myr}$; voor
$10^2$ is $\rho$ \num{1000} keer kleiner en $t_{\mathrm{ff}}$ $\sqrt{1000} = 31.6$
keer langer, \qty{3.38}{Myr}. De \omterm{def:b2:ism-star-formation:molecular-cloud}{dichte kern} stort 32 keer sneller in.
\end{solution}

\begin{solution}{exo:b2:ism-star-formation:5}
Met $k_{\mathrm B}T/(\mu m_u) = c_s^2$ is $(5k_{\mathrm B}T/G\mu m_u)^{3/2} =
5^{3/2}c_s^3G^{-3/2}$, wat de formule geeft. Omdat $M_J \propto \rho^{-1/2}$,
vergt het halveren van $M_J$ een vier keer hogere dichtheid. Omdat $M_J \propto T^{3/2}$,
vermenigvuldigt opwarmen van \qty{10}{K} tot \qty{20}{K} haar met $2^{3/2} = 2.83$.
\end{solution}

\begin{solution}{exo:b2:ism-star-formation:6}
$t_{\mathrm{KH}} = GM^2/(RL) = \num{6.674e-11}\times(\num{1.989e30})^2/(2\times
\num{6.96e8}\times2\times\num{3.83e26}) = \qty{2.48e14}{s}$, dat is
\qty{7.8}{Myr}. Dat is langer dan de levensduur van de meeste schijven (de helft is na
ongeveer \qty{3}{Myr} verdwenen): de ster trekt nog samen wanneer haar schijf verdwijnt.
\end{solution}

\begin{solution}{exo:b2:ism-star-formation:7}
Bij \qty{10}{K} is $c_s^3/G = 189^3/\num{6.674e-11} = \qty{1.01e17}{kg.s^{-1}}$, dat is
$\num{1.01e17}\times\num{3.156e7}/\num{1.989e30} = \num{1.6e-6}$ zonsmassa per
jaar, en één zonsmassa vergt \qty{0.62}{Myr}. Bij \qty{20}{K} is het tempo
$2^{3/2} = 2.83$ keer hoger, \num{4.5e-6} zonsmassa per jaar, en één zonsmassa
vergt \qty{0.22}{Myr}.
\end{solution}

\begin{solution}{exo:b2:ism-star-formation:8}
$r = \sqrt{L/(16\pi\sigma)}/T^2$ met $\sqrt{\num{3.83e26}/(16\pi\times\num{5.67e-8})}
= \qty{1.159e16}{m.K^2}$: bij \qty{1300}{K} is $r = \qty{6.86e9}{m} = \qty{0.046}{au}$,
ruim binnen de baan van Mercurius. Op \qty{5.2}{au} is $T = 278/\sqrt{5.2} = \qty{122}{K}$.
\end{solution}

\begin{solution}{exo:b2:ism-star-formation:9}
$v = \sqrt{GM/r} = \sqrt{\num{6.674e-11}\times1.3\times\num{1.989e30}/\num{1.496e13}}
= \qty{3.40}{km.s^{-1}}$, en $P = 2\pi r/v = \qty{2.77e10}{s} = 877$ jaar
(of, gelijkwaardig, $\sqrt{100^3/1.3}$ jaar). In een miljoen jaar heeft het gas
$10^6/877 \approx \num{1140}$ omlopen gemaakt: de schijf is oud vergeleken met haar eigen rotatie,
en daarom zijn haar ringen gladde cirkels.
\end{solution}

\begin{solution}{exo:b2:ism-star-formation:10}
$M = \int_0^R 4\pi r^2\rho_0R/r\,\mathrm dr = 2\pi\rho_0R^3$, en binnen $r$ is
$m(r) = 2\pi\rho_0Rr^2$. Met $\mathrm dm = 4\pi\rho_0Rr\,\mathrm dr$ is
\[
  U = -\int_0^R \frac{G\,m(r)}{r}\,\mathrm dm
    = -8\pi^2G\rho_0^2R^2\int_0^R r^2\,\mathrm dr = -\frac{8\pi^2}{3}G\rho_0^2R^5
    = -\frac23\,\frac{GM^2}{R},
\]
waarbij $\rho_0 = M/(2\pi R^3)$. Dat is $\tfrac{2/3}{3/5} = 1.11$ keer de homogene waarde:
de geconcentreerde wolk is sterker gebonden en heeft \qty{11}{\percent} meer
thermische energie nodig om gedragen te worden; ze stort gemakkelijker in.
\end{solution}

\begin{solution}{exo:b2:ism-star-formation:11}
Periode $2\pi/\Omega = \qty{6.28e14}{s}$, ongeveer \qty{20}{Myr}. Met
$R = \qty{3.086e15}{m}$ is $j = \Omega R^2 = \qty{9.52e16}{m^2.s^{-1}}$, en
$r_c = j^2/(GM) = \num{9.07e33}/\num{1.327e20} = \qty{6.83e13}{m} = \qty{457}{au}$.
Bij rotatie als een star lichaam groeit $j = \Omega r^2$ met het kwadraat van de afstand tot de as:
gas uit het binnenste van de kern heeft minder impulsmoment en komt dichter bij de ster terecht.
\end{solution}

\begin{solution}{exo:b2:ism-star-formation:12}
$\bar t = 2.46$, $\bar f = 0.588$; $\sum(t - \bar t)^2 = 12.612$,
$\sum(t - \bar t)(f - \bar f) = -2.018$, dus $b = 2.018/12.612 = 0.160$ per miljoen
jaar en $a = 0.588 + 0.160\times2.46 = 0.982$. De rechte bereikt $0.5$ bij
$(0.982 - 0.5)/0.160 = \qty{3.0}{Myr}$ en nul bij $0.982/0.160 = \qty{6.1}{Myr}$, zoals in
de figuur. Een fractie kan niet onder nul dalen: oudere sterrenhopen, zoals NGC 1960 bij
\qty{30}{Myr}, behouden een kleine fractie dicht bij nul die een rechte lijn niet kan beschrijven.
\end{solution}

\begin{solution}{pb:b2:ism-star-formation:1}
\textbf{1.} $R = 0.08\times\num{3.086e16} = \qty{2.469e15}{m}$.

\textbf{2.} $N = \qty{1.5e26}{m^{-2}}$, dus $2.8\times\num{1.661e-27}\times\num{1.5e26} =
\qty{0.698}{kg.m^{-2}}$.

\textbf{3.} $M = 0.698\times\pi R^2 = 0.698\times\num{1.915e31} = \qty{1.34e31}{kg}$,
dat is $6.7$ zonsmassa's.

\textbf{4.} $V = \tfrac43\pi R^3 = \qty{6.30e46}{m^3}$ en
$\rho = M/V = \qty{2.12e-16}{kg.m^{-3}}$.

\textbf{5.} $n = \rho/(2.33m_u) = \qty{5.48e10}{m^{-3}}$, dat is \num{5.5e4} deeltjes
per kubieke centimeter.

\textbf{6.} Ja: zijn dichtheid ligt tussen $10^4$ en $10^5$, zijn straal tussen 0.03 en
\qty{0.2}{pc}; zijn massa, 6.7 zonsmassa's, ligt aan de bovenkant van het bereik van kernen.

\textbf{7.} $M_J = 5.46\times(10^4/\num{5.48e4})^{1/2} = 2.33$ zonsmassa's.

\textbf{8.} $6.7/2.33 = 2.9$ \omterm{def:b2:ism-star-formation:jeans}{Jeans-massa}'s: de kern is instabiel en stort in.

\textbf{9.} $R_J = 0.1316\times(10^4/\num{5.48e4})^{1/2} = \qty{0.056}{pc}$, kleiner dan
de straal, \qty{0.08}{pc}: dezelfde conclusie.

\textbf{10.} $2^{3/2}\times2.33 = 6.6$ zonsmassa's: ongeveer de massa van de kern. Bij
\qty{20}{K} zou de kern op de grens liggen; het is de kou die hem doet ineenstorten.

\textbf{11.} $t_{\mathrm{ff}} = \sqrt{3\pi/(32\times\num{6.674e-11}\times\num{2.12e-16})}
= \qty{4.56e12}{s}$, ongeveer \num{145000} jaar.

\textbf{12.} $R/c_s = \num{2.469e15}/189 = \qty{1.31e13}{s}$, \num{414000} jaar: drie
keer langer dan de vrije val. Een drukgolf kan de kern niet doorkruisen voordat die
ineengestort is; de druk kan niet snel genoeg reageren om de ineenstorting te stoppen.

\textbf{13.} $c_s = \qty{189}{m.s^{-1}}$ (\cref{exo:b2:ism-star-formation:3}).

\textbf{14.} $c_s^3/G = \qty{1.01e17}{kg.s^{-1}}$, \num{1.6e-6} zonsmassa per jaar.

\textbf{15.} $1/\num{1.6e-6} = \num{624000}$ jaar.

\textbf{16.} Een derde van 6.7 is $2.2$: twee sterren van één zonsmassa.

\textbf{17.} $t_{\mathrm{KH}} = GM_\odot^2/(3R_\odot\times4L_\odot) = \qty{8.25e13}{s}$,
\qty{2.6}{Myr}.

\textbf{18.} Langer, vier keer de accretietijd: wanneer de accretie stopt, trekt de ster
nog samen en straalt ze zonder fusie; het is een \omterm{def:b2:ism-star-formation:t-tauri}{T-Tauri-ster}, nog niet op de
hoofdreeks.

\textbf{19.} $\Omega = \qty{0.4e3}{m.s^{-1}}/\qty{3.086e16}{m} = \qty{1.30e-14}{rad.s^{-1}}$.

\textbf{20.} $j = 1.30\times10^{-14}\times(\num{2.469e15})^2 = \qty{7.90e16}{m^2.s^{-1}}$.

\textbf{21.} $r_c = j^2/(GM_\odot) = \num{6.24e33}/\num{1.327e20} = \qty{4.70e13}{m}$, dat
is \qty{314}{au}: groter dan de schijf van HL Tauri, zoals te verwachten voor gas van de rand van de
kern.

\textbf{22.} $T = (L/16\pi\sigma r^2)^{1/4} = \qty{278}{K}$ op \qty{1}{au}, en
$278/\sqrt{30} = \qty{51}{K}$ op \qty{30}{au}.

\textbf{23.} $r = \num{1.159e16}/180^2 = \qty{3.58e11}{m}$, dat is \qty{2.4}{au}, zoals in de tekst.

\textbf{24.} De ster van één zonsmassa heeft
\textbf{\qty[text-series-to-math=true]{624}{} duizend jaar} nodig om zich op te bouwen, ongeveer vier
keer de \omterm{def:b2:ism-star-formation:free-fall}{vrijevaltijd} van haar kern (\num{145000} jaar): het buitenste gas bereikt de
\omterm{def:b2:ism-star-formation:protostar}{protoster} pas geleidelijk, van binnen naar buiten.
\end{solution}
